% =================================================================
% CS495 Capstone in Data Science -- Final Presentation
% Mapping Best Empirical Models: A Meta-Analysis of Winning
% Solutions in Kaggle Tabular Competitions
%
% Author: Kenneth Young
% Instructor: Pedro Albuquerque, PhD -- Bellevue College
% =================================================================

\documentclass[aspectratio=169,11pt]{beamer}

% --- Theme ---
\usetheme{metropolis}
\setbeamertemplate{frame numbering}[fraction]

% --- Packages ---
\usepackage[utf8]{inputenc}
\usepackage[T1]{fontenc}
\usepackage{graphicx}
\usepackage{booktabs}
\usepackage{array}
\usepackage{xcolor}
\usepackage{hyperref}
\usepackage{appendixnumberbeamer}

% --- Speaker notes: comment out the \setbeameroption line below
%     when you want a clean handout. Use it as-is while rehearsing. ---
% \setbeameroption{show notes on second screen=right} % uncomment for dual-monitor rehearsal
% \setbeameroption{show only notes}                    % uncomment to print notes-only deck

% --- Custom colors (lightly Bellevue-flavored: slate blue + warm accent) ---
\definecolor{kc-accent}{HTML}{C8932A}   % warm gold accent for emphasis
\definecolor{kc-muted}{HTML}{8A8A8A}    % muted gray for secondary text

% Metropolis tweaks
\setbeamercolor{progress bar}{fg=kc-accent}
\setbeamercolor{frametitle}{bg=mDarkTeal}
\setbeamercolor{alerted text}{fg=kc-accent}

% --- Title block ---
\title{Mapping Best Empirical Models}
\subtitle{A Meta-Analysis of Winning Solutions in Kaggle Tabular Competitions}
\author{Kenneth Young}
\institute{CS495 Capstone in Data Science \\ Bellevue College \\ Instructor: Pedro Albuquerque, PhD}
\date{May 2026}

% =================================================================
\begin{document}
% =================================================================

% ---------------- Title slide ----------------
\maketitle

\note{
Good [morning/afternoon]. My name is Kenneth Young, and over the next about six minutes
I'll walk you through a meta-analysis of how people win tabular machine-learning
competitions on the Kaggle platform.

The audience for this talk is mostly people who haven't competed on Kaggle, so I'll start
with what Kaggle actually is, then move into the research question, the method, and the
findings.
}


% =================================================================
\section{Introduction}
% =================================================================


% ---------------- What is Kaggle ----------------
\begin{frame}{What is Kaggle?}
\begin{itemize}
  \item Largest public platform for machine-learning \alert{competitions}
  \item Participants submit predictive models against a shared dataset
  \item Submissions are ranked by a stated metric -- accuracy, RMSE, AUC, etc.
  \item Two competition types in scope here:
  \begin{itemize}
    \item \textbf{Featured} -- monetized, sponsor-funded, real prize money
    \item \textbf{Playground Series} -- monthly, non-monetized, synthetic data
  \end{itemize}
  \item Solutions are publicly documented after the contest closes
  \begin{itemize}
    \item Writeup post (qualitative description)
    \item Notebook (executable code)
    \item Leaderboard (final rank and score)
  \end{itemize}
\end{itemize}

\note{
Kaggle is the dominant public venue for competitive machine learning. Think of it as
a series of contests: a sponsor or Kaggle itself publishes a dataset, sets an evaluation
metric, and competitors spend a month or two trying to build the best-scoring model.

There are two competition types I'm studying. Featured competitions have real prize
money -- typically tens of thousands of dollars. Playground Series competitions are
monthly contests that Kaggle runs themselves with synthetic data, primarily as a
learning venue.

After the competition ends, the winners write up what they did. That writeup is the
raw material for this study.
}
\end{frame}


% ---------------- Research question ----------------
\begin{frame}{The research question}
\begin{itemize}
  \item Thousands of writeups exist -- but no systematic empirical characterization
  \item \alert{What strategies actually win tabular Kaggle competitions?}
  \item Three concrete questions:
  \begin{enumerate}
    \item What is the \textbf{distributional structure} of winning strategies?
    \item Can the diversity be reduced to a \textbf{small typology}?
    \item Do informally observed \textbf{constraint-to-strategy patterns} hold up?
  \end{enumerate}
  \item Most existing knowledge is \textit{anecdotal} -- blog posts, individual tutorials
  \item Academic benchmarking studies compare \textit{algorithms}, not \textit{strategies}
\end{itemize}

\note{
Despite thousands of publicly documented winning solutions, nobody has systematically
characterized what those solutions look like as a population. The practitioner literature
is anecdotal -- you'll read a blog post from one prominent Kaggler about "what worked for
me." The academic literature compares algorithms like XGBoost vs LightGBM on fixed
benchmarks. Neither tells you what winning solutions across the platform actually do.

That's the gap. My three concrete questions are: what's the distribution of strategies,
can I find a small typology, and do the informal patterns hold up under quantitative
scrutiny.
}
\end{frame}


% ---------------- Why this matters ----------------
\begin{frame}{Why this matters}
\begin{columns}[T,onlytextwidth]
\column{0.55\textwidth}
\begin{itemize}
  \item \textbf{For practitioners:} replace recipe-based folk knowledge with empirically-grounded patterns
  \item \textbf{For ML meta-research:} demonstrate a replicable two-pass coding methodology
  \item \textbf{For the Kaggle community:} test whether widely-cited techniques actually dominate
\end{itemize}
\column{0.42\textwidth}
\centering
\small
\textit{``How on Earth did I win this competition?''} \\[4pt]
\textit{``1st place. That was unexpected\ldots''} \\[4pt]
\textit{``Surprisingly \#1''}
\\[10pt]
{\color{kc-muted}\scriptsize Three of five surprised-author winners in the corpus}
\end{columns}

\note{
This matters to three different audiences. For practitioners new to competitive ML,
the typology gives them an empirical map instead of recipe-based folk knowledge.
For ML meta-research, the two-pass coding method I'll describe is replicable elsewhere.
And for the Kaggle community itself, I'll show that several widely-cited patterns
don't actually dominate the data.

The pull-quotes on the right are interesting: five of the 45 winners in my corpus
explicitly described their wins as a surprise. That's not zero. There's some genuine
variance in this game that gets papered over in folk knowledge.
}
\end{frame}


% =================================================================
\section{Prior Studies}
% =================================================================


% ---------------- Related work + gap ----------------
\begin{frame}{Related work + the gap}
\begin{itemize}
  \item \textbf{Tabular ML benchmarking} -- Caruana \& Niculescu-Mizil; Erickson et al.\ (AutoGluon); Gorishniy et al.\ (numerical embeddings, NeurIPS 2022); Borisov et al.\ (tabular DL survey)
  \begin{itemize}
    \item Compares \textit{algorithms} on standardized datasets
  \end{itemize}
  \item \textbf{Competition-as-method studies} -- M5 forecasting retrospectives (Makridakis et al.); Athey et al.\ on data science competitions
  \begin{itemize}
    \item Focuses on \textit{specific competitions or families}
  \end{itemize}
  \item \textbf{Open-community attribution} -- von Krogh \& von Hippel (open source); Merton (norms of science)
  \begin{itemize}
    \item Not previously applied to competitive ML
  \end{itemize}
\end{itemize}
\vspace{8pt}
\alert{Gap:} The literature treats Kaggle as a \emph{source of techniques}, not as a \emph{community to characterize empirically}.

\note{
The existing literature breaks into three buckets. Tabular ML benchmarking compares
algorithms head-to-head on standardized datasets -- it tells you XGBoost beats logistic
regression, but not what winning solutions actually do. Competition-as-method studies
exist, but they focus on individual competitions like the M5 forecasting series or the
Netflix Prize. And there's a separate literature on attribution in open communities
from open-source software and science studies, but it hasn't been applied to competitive
ML.

The gap is straightforward: nobody's looked at recent tabular Kaggle competitions as
a population to be empirically described.
}
\end{frame}


% =================================================================
\section{Methodology}
% =================================================================


% ---------------- Data + scope ----------------
\begin{frame}{Data: 45 winning solutions, 2022--2026}
\begin{columns}[T,onlytextwidth]
\column{0.55\textwidth}
\begin{itemize}
  \item \textbf{Inclusion criteria}
  \begin{itemize}
    \item Tabular competitions only (no images / NLP / time series)
    \item At least partial writeup available
    \item Top-three finishers
  \end{itemize}
  \item \textbf{Unit of analysis}: one solution per competition (most-complete writeup of top-three)
  \item \textbf{Date range}: Feb 2022 -- Apr 2026
\end{itemize}
\column{0.42\textwidth}
\centering
\begin{tabular}{lr}
\toprule
\textbf{Cohort} & \textbf{n} \\
\midrule
PS Season 3 (2023) & 17 \\
PS Season 4 (2024) & 10 \\
PS Season 5 (2025) & 10 \\
PS Season 6 (2026) & 4 \\
TPS (2022) & 3 \\
Featured (ICR) & 1 \\
\midrule
\textbf{Total} & \textbf{45} \\
\bottomrule
\end{tabular}
\end{columns}

\note{
The corpus is 45 winning tabular Kaggle solutions from February 2022 through April 2026.
I restricted to tabular competitions, excluded time series, image and NLP work, and took
one solution per competition -- specifically, the top-three finisher with the most
complete writeup.

The cohort breakdown is roughly even across the recent Playground Series seasons, plus
three earlier TPS competitions from 2022 and one Featured competition: ICR, which was
a medical-data contest with a real prize pool.
}
\end{frame}


% ---------------- Two-pass coding ----------------
\begin{frame}{Two-pass coding design}
\begin{columns}[T,onlytextwidth]
\column{0.48\textwidth}
\textbf{Pass 1: structured schema}
\begin{itemize}
  \item 38-variable spreadsheet
  \item Models, ensembling, CV strategy, encoding, FE techniques
  \item Captures \textit{what was used}
  \item Suitable for distributional analysis
\end{itemize}

\column{0.48\textwidth}
\textbf{Pass 2: qualitative re-evaluation}
\begin{itemize}
  \item Per-writeup 9-section template
  \item Origination vs.\ inheritance
  \item Paradigm classification
  \item Captures \textit{what was originated}
  \item Suitable for typology + community-graph analysis
\end{itemize}
\end{columns}
\vspace{12pt}
\centering
\alert{Pass 1 alone cannot tell you whether a technique was originated by the winner or inherited from a publicly shared community notebook.}

\note{
The methodological centerpiece is this two-pass design. Pass 1 is a structured
spreadsheet: I went through each writeup and coded 38 variables -- what model family,
what ensembling, what cross-validation, and so on. That gives me distributional data:
how often does XGBoost appear, what's the modal CV strategy, etc.

But Pass 1 has a critical limitation: it can't tell me whether a technique listed
in the spreadsheet was actually originated by the winner, or whether they forked it
from a public notebook by another community member. That distinction matters for
the typology and especially for the community-attribution analysis. So I did a
second pass, re-reading each writeup with a nine-section template that explicitly
separates origination from inheritance.

The two-pass design is itself a contribution of the project. It's replicable.
}
\end{frame}


% ---------------- Audit + corrections ----------------
\begin{frame}{Data-quality audit + corrections workflow}
\begin{itemize}
  \item Pass 2 surfaced cell-level inconsistencies in Pass 1
  \item Built a dedicated \texttt{Data Quality Audit} sheet:
  \begin{itemize}
    \item 24 issues catalogued (5 high / 14 medium / 5 low)
    \item 17 cell-level + 7 schema-wide
  \end{itemize}
  \item Applied corrections in severity order
  \item Recorded every diff in a \texttt{Corrections Log} sheet
  \item \textbf{103 total corrections logged} (13 cell-level + 90 column-split backfills)
  \item Two new columns added: \texttt{top\_3\_margin}, \texttt{distribution\_shift\_type}
  \item One column split: \texttt{original\_data\_usage} $\to$ availability + use-mode
\end{itemize}

\note{
When I re-read the writeups for Pass 2, I started finding cell-level errors in Pass 1.
For example, one entry had n_rows coded as 300,000 when the writeup explicitly said
4 million training rows were loaded -- I had counted one CSV file but not the second.

I built a dedicated audit sheet to catalog every issue I found, colored by severity.
Then I applied the corrections in severity order -- high first, then medium, then low --
and recorded every change in a separate corrections log with the old value, the new
value, the reason, and the source. 103 total corrections.

I also added two new columns the audit identified as missing -- top-three margin
which captures how tight each competition's photo-finish was, and distribution shift
type which captures what kind of train-test shift was present. And I split one
column that was conflating two questions.

This audit trail is the methodological piece that makes the work replicable and
defensible.
}
\end{frame}


% ---------------- Limitations ----------------
\begin{frame}{Limitations}
\begin{itemize}
  \item \textbf{Small sample.} N = 45 limits sub-group analysis
  \begin{itemize}
    \item Ensemble-stacking stratum (n = 28) is reliable
    \item Smaller paradigms (n = 3--6) are descriptive only
  \end{itemize}
  \item \textbf{Single-researcher coding.} No formal inter-rater reliability
  \begin{itemize}
    \item Mitigated by Pass-2 audit on Pass-1 codings (17 corrections applied)
  \end{itemize}
  \item \textbf{Writeup-availability bias.} Excludes wins without published writeups
  \item \textbf{Winners only.} Cannot directly compare to runner-up or non-top-three solutions
  \item \textbf{Detection-limited couplings.} 200-char summary field may miss subtle signals
\end{itemize}

\note{
The limitations to be transparent about. Sample size of 45 is small for sub-group work.
The dominant paradigm has 28 entries which supports cross-tabulation reasonably well,
but the smaller paradigms only have three to six entries, so claims about them are
descriptive only.

I'm the only coder; no inter-rater reliability check. The Pass 2 audit on Pass 1
codings partially substitutes for that -- 17 cell-level corrections came out of it --
but it's not the same as a second coder.

I only see solutions with writeups, which is a real selection effect.

And there's a methodological constraint on the coupling-evidence framework I'll show
next: the strategy detector keys off a 200-character summary field, and for couplings
whose strategy signal is subtle, the actual contradicted-versus-supported split may
be partly detection-limited rather than actually contradicted.
}
\end{frame}


% =================================================================
\section{Results}
% =================================================================


% ---------------- Four paradigms ----------------
\begin{frame}{Four winning paradigms (+ two minor categories)}
\centering
\includegraphics[width=0.88\textwidth]{phase5_21_paradigm_distribution.png}
\\[6pt]
{\color{kc-muted}\scriptsize Source: \texttt{phase5\_21\_paradigm\_distribution.png}}

\note{
The first headline. The 45 winning solutions reduce to a four-paradigm typology that
captures 96 percent of the corpus.

Ensemble-stacking dominates: 28 of 45 wins, 62 percent. The winner trains many base
models -- gradient-boosted trees, neural networks, sometimes classical kernel methods --
then combines them with a stacker layer.

Single-model heavy-FE -- 13 percent. One model, hundreds of engineered features. This
is the cdeotte signature recipe.

Lookup-exploit -- 9 percent. The winner exploits a property of the synthetic data
generator to predict by direct key matching instead of by a learned model.

Problem-fit neural network -- 7 percent. Custom NN architecture matched to a specific
structural property of the data.

And two smaller categories: community-template-tweak which is fork-plus-meaningful-tweak,
and mixed for one entry that doesn't fit cleanly.

Important descriptive finding: there are zero entries scored as pure forks. Even the
most fork-heavy paradigm contains at least one meaningful original contribution per
winner.
}
\end{frame}


% ---------------- Coupling evidence (the headline) ----------------
\begin{frame}{Coupling evidence: only 2 of 6 promoted patterns survive}
\centering
\includegraphics[width=0.92\textwidth]{phase5_51_coupling_evidence.png}
\\[6pt]
{\color{kc-muted}\scriptsize Source: \texttt{phase5\_51\_coupling\_evidence.png}}

\note{
This is the analytical centerpiece. Pass 2 identified six promoted constraint-to-strategy
couplings -- patterns that appeared in multiple writeups and looked like real community
practice. To test each one quantitatively, I defined a constraint filter and a strategy
filter, then counted supporting, contradicting, and neutral cases.

Two couplings survive at the strong-evidence threshold of 75 percent support.

Linear-stacker dominance: 32 of 40 testable cases, 80 percent. When winners ensemble,
they overwhelmingly choose a linear stacker -- Ridge regression, mean blending, hill
climbing -- rather than a nonlinear one.

Distribution-shift-driven custom CV: four of five testable cases, also 80 percent.
But the testable sample is only five entries, so this is a strong-with-small-n caveat.

Four couplings are contradicted. The most striking is "use available external original
as columns" -- the widely-cited cdeotte pattern. Of 30 entries with an available
external original, only 11 use it as columns or features-derived. The dominant pattern
is actually plain rows-only concatenation. The "use as columns" framing is one author's
practice, not the community's practice.

The two contradicted couplings labeled in red and the photo-finish coupling are partly
detection-limited, which I flagged in the limitations.
}
\end{frame}


% ---------------- Paradigm by era ----------------
\begin{frame}{Paradigm composition by era}
\centering
\includegraphics[width=0.88\textwidth]{phase5_22_paradigm_by_era.png}
\\[6pt]
{\color{kc-muted}\scriptsize Source: \texttt{phase5\_22\_paradigm\_by\_era.png}}

\note{
Paradigm composition is era-dependent. The three earliest competitions in the corpus --
the Tabular Playground Series from 2022 -- contributed two problem-fit-NN wins and one
lookup-exploit win, and zero ensemble-stacking wins.

From Season 3 onward, ensemble-stacking is the modal paradigm in every era, accounting
for 60 to 80 percent of each era's wins.

The diversity of paradigms peaks in Season 3, where five of the six paradigm categories
appear at least once.

This era-dependence matters: the typology is current as of mid-2026, but earlier eras
of competitive tabular ML looked qualitatively different. The dataset captures both
the current dominant pattern and the historical variation.
}
\end{frame}


% ---------------- Community + cdeotte timeline ----------------
\begin{frame}{Community centrality: 3 archetypes}
\centering
\includegraphics[width=0.95\textwidth]{phase5_41_author_centrality.png}
\\[4pt]
{\color{kc-muted}\scriptsize Source: \texttt{phase5\_41\_author\_centrality.png}}

\note{
Aggregating wins and citations across the corpus, three centrality archetypes appear.

First: pure technique-source. The most-cited handle in the entire corpus is siukeitin,
with 12 citations across 12 different writeups. siukeitin never appears as a top-three
winner in any of my 45 entries. paddykb is similar: 5 citations, zero wins. These
are community members whose techniques propagate widely but who don't themselves win.

Second: member who occasionally wins. tilii7 has one win plus eight citations.
masayakawamata has one win plus five citations.

Third: dominant competitor who is also cited. cdeotte has four wins plus six citations.
Same for mahoganybuttstrings and for the AmbrosM-plus-team-partner pair.

This empirical separation -- pure source, occasional winner, dominant cited competitor --
is the strongest example I have of the originator-versus-canonizer gap in competitive ML.
It's worth saying out loud: the most-cited technique source in my corpus has not won a
top-three position in any of the competitions I analyzed.
}
\end{frame}


% ---------------- cdeotte timeline ----------------
\begin{frame}{Observer-to-winner: 22 months}
\centering
\includegraphics[width=0.95\textwidth]{phase5_42_cdeotte_timeline.png}
\\[6pt]
{\color{kc-muted}\scriptsize Source: \texttt{phase5\_42\_cdeotte\_timeline.png}}

\note{
Zooming in on one of the most prominent winners: cdeotte. His earliest documented
presence in my corpus is February 2023, when he commented on a winning writeup as
an outside observer. His first win in the corpus is December 2024, 22 months later.

Between those two dates he appears as a commenter on other writeups, as a non-winning
competitor placing 81st in one competition and 806th in another, and as a cited
source for techniques other winners use. Then his wins accelerate, producing four
top-three finishes in 14 months.

This 22-month observe-before-win arc is a specific case of a more general pattern:
dominance in competitive ML is preceded by a long period of community observation.
The techniques that cdeotte later canonized in his writeups -- hill climbing, brute-force
feature engineering, RAPIDS XGBoost -- were already in use by other community members
before he started winning with them.
}
\end{frame}


% ---------------- Unexpected findings ----------------
\begin{frame}{Three unexpected findings}
\begin{itemize}
  \item \textbf{No pure forks among winners.} Zero entries at origination-score 0
  \begin{itemize}
    \item Even community-template-tweak winners add a meaningful tweak
  \end{itemize}
  \item \textbf{Most-cited contributor has never won.} siukeitin: 12 citations, 0 wins in corpus
  \begin{itemize}
    \item Pure technique-source role is real and measurable
  \end{itemize}
  \item \textbf{``Use as columns'' is contradicted.} 19 of 30 entries with available original use rows-only
  \begin{itemize}
    \item The widely-cited pattern reflects one author's practice, not the community's
  \end{itemize}
\end{itemize}

\note{
To close: three findings I didn't expect going in.

One: no pure forks among winners. I had assumed -- as a lot of community folk knowledge
implies -- that some wins are essentially "fork a strong public notebook and submit." My
data says no. Every single one of the 45 winners contributes at least one meaningful
original element.

Two: the most-cited technique source in the corpus has never won a top-three competition
in the corpus. That's siukeitin, cited 12 times across 12 different writeups. The
originator-versus-canonizer gap is not just a theoretical possibility -- it's directly
measurable here.

Three: the "use available original as columns" pattern, which the practitioner community
widely associates with one prominent author, is actually contradicted by the data. Most
winners with an available external original just concatenate the original as additional
training rows. The "as columns" pattern is one author's idiosyncratic practice that
got popularized in his writeups.

These three findings together support the broader claim that empirically characterizing
winning Kaggle practice surfaces real contradictions with folk knowledge.

I'm happy to take questions.
}
\end{frame}


% ---------------- Thank you / Q&A ----------------
\begin{frame}[standout]
Thank you. \\[12pt]
{\Large Questions?}
\end{frame}


% =================================================================
% Appendix slides -- for question time
% =================================================================
\appendix


\begin{frame}{Appendix: coupling evidence detail}
\centering
\small
\begin{tabular}{lrrrrl}
\toprule
\textbf{Coupling} & \textbf{Supp.} & \textbf{Contr.} & \textbf{n test.} & \textbf{Rate} & \textbf{Verdict} \\
\midrule
Linear stacker dominance & 32 & 8 & 40 & 80\% & Strong \\
Dist.\ shift $\to$ custom CV & 4 & 1 & 5 & 80\% & Strong (small n) \\
Heavy citations $\to$ fork-heavy & 3 & 5 & 8 & 38\% & Contradicted \\
Available original $\to$ as cols & 11 & 19 & 30 & 37\% & Contradicted \\
Small data $\to$ no-FE single & 3 & 10 & 13 & 23\% & Contradicted \\
Photo-finish $\to$ valid.\ disc. & 6 & 21 & 27 & 22\% & Contradicted$^\ast$ \\
\bottomrule
\end{tabular}
\\[8pt]
{\scriptsize $^\ast$ Detection-limited: strategy filter keys off 200-char summary; full re-eval may shift verdict.}
\end{frame}


\begin{frame}{Appendix: dataset by size and paradigm}
\centering
\includegraphics[width=0.9\textwidth]{phase5_24_paradigm_n_rows.png}
\end{frame}


\begin{frame}{Appendix: photo-finish rate by paradigm}
\centering
\includegraphics[width=0.9\textwidth]{phase5_23_paradigm_photofinish.png}
\end{frame}


\begin{frame}{Appendix: external original $\times$ paradigm}
\centering
\includegraphics[width=0.9\textwidth]{phase5_31_use_mode_paradigm.png}
\end{frame}


\end{document}
